% ============================================================================
% Section 10: Experimental Protocol
% ============================================================================

\section{Experimental protocol}
\label{sec:experimental}

Validating UHM requires a systematic, multi-phase experimental programme organized by
increasing cost, complexity, and risk. The protocol tests all 23 predictions across four
phases, beginning with experiments that require nothing more than a GPU and ending with
large-scale clinical studies. The design principle is \emph{maximal risk first}: the most
falsifiable predictions are tested earliest, so that the theory is either killed quickly
or earns the right to expensive experiments.

\subsection{Phase I: Digital validation (0--6 months)}
\label{subsec:phase1}

Phase~I requires no laboratory equipment, no human subjects, and no ethics approval.
Twelve of the 23 predictions are testable in silico on any implementation of a
\textbf{$\Gam$-native agent} --- a system whose evolution is governed by Lindbladian
dynamics $\Lind_\Omega = \Lind_0 + \Regen$ on the coherence matrix
$\Gam \in \Dens(\mathbb{C}^7)$.

\paragraph{Requirements for a $\Gam$-native agent.}
\begin{enumerate}[leftmargin=2em,nosep]
  \item \textbf{CPTP dynamics:} Evolution of $\Gam$ via a CPTP channel (T-62).
  \item \textbf{7D structure:} State space is $\Dens(\mathbb{C}^7)$ with dimensions
        [A, S, D, L, E, O, U].
  \item \textbf{Consciousness verifier:} At each step, $P$, $R$, $\Phi$, $\CohE$,
        $\sigma_k$ are computed.
  \item \textbf{Multi-phase training with hard gates:} Initialization ($P > P_{\min}$),
        foundation ($P \in (2/7, 3/7]$, $R \geq 1/3$), autonomous learning
        ($\sigma$-directed data selection).
  \item \textbf{Checkpoint system:} Full $\Gam$ state saved for perturbation tests.
  \item \textbf{GPU acceleration:} $\geq 1$ GPU with $\geq 40$\,GB for Monte Carlo
        (Experiment~I.8).
\end{enumerate}

\paragraph{Experiments.}
Table~\ref{tab:phase1} lists the 11 Phase~I experiments, each linked to a specific
prediction.

\begin{table}[H]
\centering
\caption{Phase~I experiments (digital, in silico).}
\label{tab:phase1}
\small
\begin{tabular}{@{}clll@{}}
\toprule
\textbf{Exp.} & \textbf{Prediction tested} & \textbf{Key metric} & \textbf{Falsification} \\
\midrule
I.1 & Pred 1 (No-Zombie) & $\tau_{\mathrm{death}}$ after E-suppression & $\tau_E \geq \tau_A$ \\
I.2 & Pred 7 (Stability radius) & $h_{\mathrm{crit}}^2$ vs.\ $P - 2/7$ & $R^2 < 0.9$ \\
I.3 & Pred 8 (Info capacity) & $I(\mathrm{obs};\delta\Gam)$ & $I > 2.81$ bits \\
I.4 & Pred 10 ($N=7$ learning) & Accuracy at $N=5$ & $> 75\%$ at $N=5$ \\
I.5 & Pred 12 (SAD ceiling) & $\mathrm{SAD}(\Gam)$ via T-136 & $\exists\,\Gam:\mathrm{SAD}(\Gam)\geq 4$ \\
I.6 & Pred 13 (Genesis time) & Steps to $P > 2/7$ & $n > n_{\mathrm{genesis}}$ \\
I.7 & Pred 14 (Phase coherence) & $\Phi$ with/without co-rotation & $\Phi \geq 1$ without \\
I.8 & Pred 17 (Exponents, prelim.) & $\beta$ from Monte Carlo & $\beta \notin [0.20, 0.30]$ \\
I.9 & Pred 19 (CPTP anchor) & $\|\pi - \pi_{\mathrm{can}}\|_\diamond$ & $> 0.1$ at convergence \\
I.10 & Pred 9 (Learning speed) & $n$ to 90\% accuracy & $n < n_{\mathrm{info}}$ \\
I.11 & Pred 11 ($N=7$ social) & Coordination at $N=5$ & ToM+ISL+Nash at $N=5$ \\
\bottomrule
\end{tabular}
\end{table}

\paragraph{Go/no-go criterion.}
All 11 Phase~I experiments must be confirmed before proceeding to Phase~II. Any L1 or L2
falsification halts the programme pending theoretical revision.

\subsection{Phase II: Neurocalibration (6--18 months)}
\label{subsec:phase2}

The central task of Phase~II is to build the bridge
$\pi_{\mathrm{bio}}: (\mathrm{EEG}, \mathrm{fMRI}, \mathrm{HRV}) \to \Gam \in
\Dens(\mathbb{C}^7)$ and verify that $\Pcrit = 2/7$ coincides with the empirical PCI
threshold $\mathrm{PCI}^* = 0.31$~\cite{casali2013}.

\paragraph{Formal definition of $\pi_{\mathrm{bio}}$.}
The neural anchor map is a CPTP channel
$\pi_{\mathrm{bio}}: \mathcal{O}_{\mathrm{neural}} \to \Dens(\mathbb{C}^7)$
constructed in four steps:
(1)~extract 7 canonical neural observables (spectral edge, DFA exponent, permutation
entropy, theta-gamma coupling, PCI, HRV, mean PLI) corresponding to the 7 UHM
dimensions $\{A,S,D,L,E,O,U\}$;
(2)~normalise to a density matrix diagonal via softmax:
$\gamma_{kk} = \exp(\beta z_k)/\sum_j \exp(\beta z_j)$, where $z_k = (x_k - \mu_k)/\sigma_k$
is standardised and $\beta$ is calibrated to match PCI$^* \leftrightarrow \Pcrit$;
(3)~reconstruct off-diagonal coherences from pairwise coherence and phase-locking values,
regularised by Cholesky decomposition to guarantee positivity;
(4)~align to canonical gauge via $G_2$-covariant Procrustes (T-123~\status{T}).
The construction is CPTP by composition (softmax $+$ Cholesky $=$ positive, trace-preserving).
The calibration parameter $\beta$ is empirical~[H]; after Phase~II validation, the full
map becomes a validated instrument~[T].

\paragraph{Statistical layer: the $\Gam$-tomography estimator.}
Everything downstream of the embedding is standard estimator theory rather
than an unaccountable fit. With windowed embeddings $x_n = \pi(\text{window}_n)
\in \mathbb{C}^7$, the coherence matrix is the normalized second moment
$\Gam = M/\Tr M$, $M = \mathbb{E}[x x^{\dagger}]$, and the estimator
$\widehat{\Gam}_N = \widehat{\Sigma}_N / \Tr\widehat{\Sigma}_N$ with
$\widehat{\Sigma}_N = \tfrac{1}{N}\sum_n x_n x_n^{\dagger}$ is Hermitian,
PSD, and unit-trace by construction. It is (i)~strongly consistent; (ii)~its
error obeys a matrix-Bernstein bound, $N \gtrsim C B^2/(\tau^2\varepsilon^2)
\ln(14/\delta)$ samples suffice for $\|\widehat{\Gam}_N - \Gam\|_{\mathrm{op}}
\leq \varepsilon$ with probability $1-\delta$ (rate $O(N^{-1/2})$);
(iii)~purity is estimated by the two-sample $U$-statistic
$\widehat{P}_N = \sum_{m \neq n} |x_m^{\dagger} x_n|^2 / [N(N-1)
(\Tr\widehat{\Sigma}_N)^2]$, whose numerator is exactly unbiased for
$\Tr(M^2)$ (the normalized ratio keeps only a small residual $O(1/N)$
denominator bias, opposite in sign and $\sim 3\times$ smaller than the
na\"ive plug-in $\Tr(\widehat{\Gam}_N^2)$); and (iv)~for autocorrelated
recordings all confidence statements use the effective sample size
$N_{\mathrm{eff}} = N/(2\tau_{\mathrm{corr}} + 1)$, with the integrated
autocorrelation time $\tau_{\mathrm{corr}}$ estimated by batch means. A
self-tested reference implementation (\texttt{gamma\_tomography.py},
$\sim$200 lines of NumPy) accompanies the project and reproduces the
$O(N^{-1/2})$ rate, the $U$-statistic debiasing, and the $N_{\mathrm{eff}}$
correction numerically. The single irreducible modelling choice is the
7-channel embedding $\pi$ itself; the estimation below it is rigorous.

\paragraph{The information price of self-tomography (candidate).}
The estimator above has a universal-coding twin: the cumulative
prequential cost of \emph{learning} an unknown $\Gam \in D(\mathbb{C}^d)$
from $n$ samples, over and above the cost of knowing it, is
$\tfrac{d^2-1}{2}\log_2 n + C_Q(d) + o(1)$ --- the dimension term is
fixed by the $48 = d^2{-}1$ readout theorem, and the constant admits an
exact identity. With the Bures metric normalized as one quarter of the
quantum (SLD) Fisher metric --- this paper's standing convention --- the
Jeffreys-type volume of state space obeys
$2^{\,d^2-1}\,V_{\mathrm{Bures}}(d) = \pi^{d^2/2}/\Gamma(d^2/2)$
exactly (the dyadic factor of the quarter-Fisher convention cancels the
$2^{\,1-d^2}$ of the Sommers--\.Zyczkowski volume identically), so the
Clarke--Barron-type constant \emph{collapses onto the classical family}:
$C_Q(d) = C(d^2)$, where $C(m) = \tfrac{m}{2}\log_2\pi -
\tfrac{m-1}{2}\log_2(2\pi) - \log_2\Gamma(m/2)$ is the startup constant
of a classical $m$-outcome source under the Krichevsky--Trofimov
mixture. Informationally, a quantum state of dimension $d$ begins its
life as a classical die with $d^2$ faces; for the organism's
$\Gam \in D(\mathbb{C}^7)$ the price is $24\log_2 n + C(49)$ with
$C(49) = -99.9119$~bits. Status is kept honest by layer: the volume
identity and the family collapse are exact algebra (verified to
$10^{-12}$ for $d = 2..7$ by \texttt{natal\_startup\_verify.py}, 4/4,
including the closed qubit volume $V_{\mathrm{Bures}}(2) = \pi^2/8$);
the $\tfrac{d^2-1}{2}\log_2 n$ dimension term is standard for
mixed-state universal coding; the \emph{operational} attainability of
the constant by adaptive tomography is the candidate's named open
court.

\paragraph{Equipment.}
All required instruments are commercially available and standard in
neuroscience labs (Table~\ref{tab:equipment}). No custom hardware is needed.

\begin{table}[H]
\centering
\caption{Equipment classes required for Phase~II neurocalibration.}
\label{tab:equipment}
\begin{tabular}{@{}ll@{}}
\toprule
\textbf{Component} & \textbf{Purpose} \\
\midrule
TMS-EEG (e.g.\ navigated TMS + 60-ch EEG) & Causal perturbation and evoked EEG \\
High-density EEG ($\geq 128$ channels)    & Spectral and topographic analysis \\
3\,T fMRI (standard research access)      & Spatial localisation \\
HRV / autonomic monitoring                & Peripheral physiological covariates \\
Polysomnography                           & Sleep-stage classification \\
MRI-compatible neuronavigation            & TMS targeting accuracy \\
\bottomrule
\end{tabular}
\end{table}

\paragraph{Key experiments.}
\begin{itemize}[leftmargin=2em]
  \item \textbf{Exp.~II.1: $\Pcrit \leftrightarrow \mathrm{PCI}^*$} (the most important
        experiment). Propofol-induced loss of consciousness with TMS-EEG, $N = 50$ subjects.
        At five propofol target levels ($C_e = 0.5, 1.0, 1.5, 2.0, 2.5$\,$\mu$g/ml), apply
        $\pi_{\mathrm{bio}}$ and compute $P$. Primary outcome:
        $P_{\mathrm{boundary}} = 2/7 \pm 0.05$ (95\% CI includes $2/7$). Statistical plan:
        one-sample $t$-test, $\alpha = 0.01$, power $> 0.95$ at $\mathrm{SD} = 0.06$.
        Falsification: $|P_{\mathrm{boundary}} - 2/7| > 0.1$ at $p < 0.01$.
  \item \textbf{Exp.~II.2: Critical exponents} (the riskiest experiment). TMS-EEG during
        full-night polysomnography (W$\to$N1$\to$N2$\to$N3$\to$REM$\to$W), $N = 50$
        subjects, 100 TMS trials every 15 minutes, yielding $\sim$1600 data points. For
        each: compute PCI and $P$. Fit $\mathrm{PCI} \sim (P - \Pcrit)^\beta$. Prediction:
        $\beta = 1/4 \pm 0.05$. Additional: $\nu = 1/2$ from spatial extent of the TMS-evoked
        response; $\gamma = 1$ from amplitude of PCI variability near threshold. Analysis:
        nonlinear regression, bootstrap 95\% CI, comparison with ordinary mean-field ($\beta = 1/2$)
        and 3D~Ising ($\beta \approx 0.326$). Falsification: $\beta \notin [0.20, 0.30]$
        at $p < 0.01$.
  \item \textbf{Exp.~II.3: Ignition dynamics} (Pred~16). Measure latency $T_{\mathrm{ign}}$
        from TMS pulse to first PCI burst ($> 50\%$ of $\mathrm{PCI}_{\mathrm{wake}}$),
        $N = 30$ (subsample of Exp.~II.1). Prediction:
        $T_{\mathrm{ign}} \sim (P - \Pcrit)^{-1}\kappa_0^{-1}$ (divergence near threshold).
        Falsification: $T_{\mathrm{ign}}$ does not depend on $(P - \Pcrit)$, $R^2 < 0.3$.
  \item \textbf{Exp.~II.4: Spectral gap and gamma rhythm} (Pred~22~\status{H}). HD-EEG 128-ch,
        resting state (10 min, eyes open and closed), $N = 30$. Spectral analysis: dominant
        frequency in gamma range (30--100\,Hz). Calibrate $\lambda_{\mathrm{gap}}$ from
        Lindbladian parameters. Prediction:
        $\nu_{\mathrm{predicted}} = \lambda_{\mathrm{gap}}/(2\pi)$ matches gamma range.
        Falsification: $\lambda_{\mathrm{gap}}/(2\pi) \notin [10, 200]$\,Hz.
\end{itemize}

\subsection{Phase III: Clinical validation (18--36 months)}
\label{subsec:phase3}

Phase~III applies UHM to disorders of consciousness (DOC), where the theory's predictions
have the most direct clinical significance.

\begin{itemize}[leftmargin=2em]
  \item \textbf{Exp.~III.1: DOC classification} (Pred~21). $N = 80$ (20 coma, 20 minimally
        conscious state [MCS], 20 vegetative state [VS/UWS], 20 healthy controls). TMS-EEG
        + fMRI + HRV $\to \pi_{\mathrm{bio}} \to \Gam$. Prediction:
        $P(\Gam_{\mathrm{MCS}}) > 2/7$ for $\geq 90\%$ of MCS patients;
        $P(\Gam_{\mathrm{VS}}) < 2/7$ for $\geq 80\%$. Falsification: sensitivity $< 80\%$
        or specificity $< 75\%$.
  \item \textbf{Exp.~III.2: E-coherence and recovery} (Pred~2). $N = 60$, stroke
        rehabilitation. Measure $\CohE$ at admission via EEG $\to \pi_{\mathrm{bio}}$.
        Assess recovery at 3 months (Barthel Index, modified Rankin Scale). Prediction:
        $r > 0.3$ (Pearson) between $\CohE(t_0)$ and recovery rate. Control for age and
        physical health. Falsification: $r \leq 0$ at $N = 60$, $p < 0.05$.
  \item \textbf{Exp.~III.3: Resting-state attractor} (Pred~15). $N = 30$, healthy,
        5 sessions on different days. EEG + fMRI $\to \pi_{\mathrm{bio}} \to P$. Prediction:
        $P_{\mathrm{rest}} \to 3/7 \pm 0.05$ (T-124). Falsification:
        $|P_{\mathrm{mean}} - 3/7| > 0.1$ at $N = 30$.
\end{itemize}

\subsection{Phase IV: Cognitive and social validation (12--24 months)}
\label{subsec:phase4}

Phase~IV tests UHM's predictions about the structure of cognition, collective consciousness,
and language independence.

\begin{itemize}[leftmargin=2em]
  \item \textbf{Exp.~IV.1: 7D stress tensor} (Pred~3). Compile a database of 200+
        stressors from psychology, medicine, and organizational science. Five independent
        experts classify each stressor by 7 components [A, S, D, L, E, O, U]. Inter-rater
        reliability: Cohen's $\kappa$. Prediction: 100\% coverage (every stressor maps to
        $\geq 1$ component); residual category is empty. Falsification: a stressor
        unclassifiable by $\geq 4/5$ experts.
  \item \textbf{Exp.~IV.2: Collective consciousness} (Pred~5). Hyperscanning EEG
        ($4 \times 32$-ch, synchronization via LSL). Ten jazz quartets (coordinated) vs.\
        ten random musician groups (uncoordinated). Compute $\Phi_\otimes$ for each group.
        Falsification: $\Phi_\otimes(\mathrm{coordinated}) \leq
        \Phi_\otimes(\mathrm{uncoordinated})$ at $p < 0.05$ (Mann--Whitney).
  \item \textbf{Exp.~IV.3: Pre-linguistic cognition} (Pred~4). $N = 30$ (15 Broca's aphasia,
        15 healthy controls). Battery of nonverbal cognitive tests: K1 (perception),
        K2 (emotions), K3 (categorization), K4 (planning). Prediction: K1--K4 scores in
        aphasic patients preserved at $> 80\%$ of normal (T-100). Falsification: K3 or K4
        decline $> 50\%$ in aphasia.
\end{itemize}

\subsection{Timeline and dependency summary}

\begin{table}[H]
\centering
\caption{Four-phase timeline and dependencies.}
\label{tab:timeline}
\begin{tabular}{@{}llll@{}}
\toprule
\textbf{Phase} & \textbf{Start} & \textbf{Duration} & \textbf{Prerequisite} \\
\midrule
I. Digital   & Month 0  & 6 mo.  & GPU ($\geq 40$\,GB) \\
II. Neuro    & Month 6  & 12 mo. & Phase I pass; IRB approval \\
III. Clinical & Month 12 & 24 mo. & Phase II $\pi_{\mathrm{bio}}$ calibration \\
IV. Cognitive & Month 12 & 12 mo. & Phase I pass \\
\bottomrule
\end{tabular}
\end{table}

\noindent
The complete programme spans approximately 36 months. The design places the
maximally falsifiable tests in Phase~I, which requires only standard GPU
hardware and no human subjects: twelve of the 23 predictions can be tested
in silico before any laboratory experiment begins. If Phase~I fails (any
L1 or L2 falsification), the programme halts immediately and no further
resources are committed. The later phases reuse equipment that is already
standard in neuroscience laboratories (TMS-EEG, HD-EEG, 3\,T fMRI); no
bespoke apparatus is required.
